\documentclass[a4paper,12pt]{article}
\usepackage[T1]{fontenc}
\usepackage[german]{babel}

\begin{document}
\thispagestyle{empty}
\begin{center}
\Large\bfseries
Ueber die Hypothesen, welche der Geometrie zu Grunde liegen.\\[12 pt]
Bernhard Riemann\\[12 pt]
[Aus dem dreizehnten Bande der Abhandlungen der K\"{o}niglichen
Gesellschaft der Wissenschaften zu G\"{o}ttingen.]\\[24 pt]
\large\mdseries
Transcribed by D. R. Wilkins\\[12 pt]
Preliminary Version: December 1998\\
Corrected: April 2000
\end{center}

\newpage
\setcounter{page}{1}


\title{Ueber die Hypothesen, welche der Geometrie zu Grunde liegen.}
\author{Bernhard Riemann}
\date{[Aus dem dreizehnten Bande der Abhandlungen der K\"{o}niglichen
Gesellschaft der Wissenschaften zu G\"{o}ttingen.]\thanks{Diese
Abhandlung ist am 10.~Juni 1854 von dem Verfasser bei dem zum Zweck
seiner Habilitation veranstalteten Colloquium mit der philosophischen
Facult\"{a}t zu G\"{o}ttingen vorgelesen worden.  Hieraus erkl\"{a}rt
sich die Form der Darstellung, in welcher die analytischen
Untersuchungen nur angedeutet werden konnten; einige Ausf\"{u}hrungen
derselben findet man in der Beantwortung der Pariser Preisaufgabe
nebst den Anmerkungen zu derselben.}}
\maketitle

\begin{center}
\large\bfseries
Plan der Untersuchung
\end{center}

Bekanntlich setzt die Geometrie sowohl den Begriff des Raumes,
als die ersten Grundbegriffe f\"{u}r die Constructionen in Raume
als etwas Gegebenes voraus.  Sie giebt von ihnen nur
Nominaldefinitionen, w\"{a}hrend die wesentlichen Bestimmungen in
Form von Axiomen auftreten.  Das Verh\"{a}ltniss dieser
Voraussetzungen bleibt dabei in Dunkeln; man sieht weder ein, ob
und in wie weit ihre Verbindung nothwendig, noch a priori, ob sie
m\"{o}glich ist.

Diese Dunkelheit wurde auch von \emph{Euklid} bis auf
\emph{Legendre}, um den ber\"{u}hmtesten neueren Bearbeiter der
Geometrie zu nennen, weder von den Mathematikern, noch von den
Philosophen, welche sich damit besch\"{a}ftigten, gehoben.  Es
hatte dies seinen Grund wohl darin, dass der allgemeine Begriff
mehrfach ausgedehnter Gr\"{o}ssen, unter welchem die
Raumgr\"{o}ssen enthalten sind, ganz unbearbeitet blieb.  Ich
habe mir daher zun\"{a}chst die Aufgabe gestellt, den Begriff
einer mehrfach ausgedehnten Gr\"{o}sse aus allgemeinen
Gr\"{o}ssenbegriffen zu construiren.  Es wird daraus hervorgehen,
dass eine mehrfach ausgedehnte Gr\"{o}sse verschiedener
Massverh\"{a}ltnisse f\"{a}hig ist und der Raum also nur einen
besonderen Fall einer dreifach ausgedehnten Gr\"{o}sse bildet.
Hiervon aber ist eine nothwendige Folge, dass die S\"{a}tze der
Geometrie sich nicht aus allgemeinen Gr\"{o}ssenbegriffen
ableiten lassen, sondern dass diejenigen Eigenschaften, durch
welche sich der Raum von anderen denkbaren dreifach ausgedehnten
Gr\"{o}ssen unterscheidet, nur aus der Erfahrung entnommen werden
k\"{o}nnen.  Hieraus entsteht die Aufgabe, die einfachsten
Thatsachen aufzusuchen, aus denen sich die Massverh\"{a}ltnisse
des Raumes bestimmen lassen---eine Aufgabe, die der Natur der
Sache nach nicht v\"{o}llig bestimmt ist; denn es lassen sich
mehrere Systeme einfacher Thatsachen angeben, welche zur
Bestimmung der Massverh\"{a}ltnisse des Raumes hinreichen; am
wichtigsten ist f\"{u}r den gegenw\"{a}rtigen Zweck das von
\emph{Euklid} zu Grunde gelegte.  Diese Thatsachen sind wie
alle Thatsachen nicht nothwendig, sondern nur von empirischer
Gewissheit, sie sind Hypothesen; man kann also ihre
Wahrscheinlichkeit, welche innerhalb der Grenzen der Beobachtung
allerdings sehr gross ist, untersuchen und hienach \"{u}ber die
Zul\"{a}ssigkeit ihrer Ausdehnung jenseits der Grenzen der
Beobachtung, sowohl nach der Seite des Unmessbargrossen, als nach
der Seite des Unmessbarkleinen urtheilen.

\begin{center}
\large\bfseries
I. Begriff einer $n \,$\-fach ausgedehnten Gr\"{o}sse.
\end{center}

Indem ich nun von diesen Aufgaben zun\"{a}chst die erste, die
Entwicklung des Begriffs mehrfach ausgedehnter Gr\"{o}ssen, zu
l\"{o}sen versuche, glaube ich um so mehr auf eine nachsichtige
Beurtheilung Anspruch machen zu d\"{u}rfen, da ich in dergleichen
Arbeiten philosophischer Natur, wo die Schwierigkeiten mehr in
den Begriffen, als in der Construction liegen, wenig ge\"{u}bt
bin und ich ausser einigen ganz kurzen Andeutungen, welche Herr
Geheimer Hofrath \emph{Gauss} in der zweiten Abhandlung
\"{u}ber die biquadratischen Reste, in den G\"{o}ttingenschen
gelehrten Anzeigen und in seiner Jubil\"{a}umsschrift dar\"{u}ber
gegeben hat, und einigen philosophischen Untersuchungen
\emph{Herbart}'s, durchaus keine Vorarbeiten benutzen konnte.

\medbreak

\centerline{1.}

\nobreak\medskip

Gr\"{o}ssenbegriffe sind nur da m\"{o}glich, wo sich ein
allgemeiner Begriff vorfindet, der verschiedene Bestimmungsweisen
zul\"{a}sst.  Je nachdem unter diesen Bestimmungsweisen von einer
zu einer andern ein stetiger Uebergang stattfindet oder nicht,
bilden sie eine stetige oder discrete Mannigfaltigkeit; die
einzelnen Bestimmungsweisen heissen im erstern Falle Punkte, im
letztern Elemente dieser Mannigfaltigkeit.  Begriffe, deren
Bestimmungsweisen eine discrete Mannigfaltigkeit bilden, sind so
h\"{a}ufig, dass sich f\"{u}r beliebig gegebene Dinge wenigstens
in den gebildeteren Sprachen immer ein Begriff auffinden
l\"{a}sst, unter welchem sie enthalten sind (und die Mathematiker
konnten daher in der Lehre von den discreten Gr\"{o}ssen
unbedenklich von der Forderung ausgehen, gegebene Dinge als
gleichartig zu betrachten), dagegen sind die Veranlassungen zur
Bildung von Begriffen, deren Bestimmungsweisen eine stetige
Mannigfaltigkeit bilden, im gemeinen Leben so selten, dass die
Orte der Sinnengegenst\"{a}nde und die Farben wohl die einzigen
einfachen Begriffe sind, deren Bestimmungsweisen eine mehrfach
ausgedehnte Mannigfaltigkeit bilden.  H\"{a}ufigere Veranlassung
zur Erzeugung und Ausbildung dieser Begriffe findet sich erst in
der h\"{o}hern Mathematik.

Bestimmte, durch ein Merkmal oder eine Grenze unterschiedene
Theile einer Mannigfaltigkeit heissen Quanta.  Ihre Vergleichung
der Quantit\"{a}t nach geschieht bei den discreten Gr\"{o}ssen
durch Z\"{a}hlung, bei den stetigen durch Messung.  Das Messen
besteht in einem Aufeinanderlegen der zu vergleichenden
Gr\"{o}ssen; zum Messen wird also ein Mittel erfordert, die eine
Gr\"{o}sse als Massstab f\"{u}r die andere fortzutragen.  Fehlt
dieses, so kann man zwei Gr\"{o}ssen nur vergleichen, wenn die
eine ein Theil der andern ist, und auch dann nur das Mehr oder
Minder, nicht das Wieviel entscheiden.  Die Untersuchungen,
welche sich in diesem Falle \"{u}ber sie anstellen lassen, bilden
einen allgemeinen von Massbestimmungen unabh\"{a}ngigen Theil der
Gr\"{o}ssenlehre, wo die Gr\"{o}ssen nicht als unabh\"{a}ngig von
der Lage existirend und nicht als durch eine Einheit
ausdr\"{u}ckbar, sondern als Gebiete in einer Mannigfaltigkeit
betrachtet werden.  Solche Untersuchungen sind f\"{u}r mehrere
Theile der Mathematik, namentlich f\"{u}r die Behandlung der
mehrwerthigen analytischen Functionen ein Bed\"{u}rfniss
geworden, und der Mangel derselben ist wohl eine Hauptursache,
dass der ber\"{u}hmte \emph{Abel}'sche Satz und die Leistungen
von \emph{Lagrange}, \emph{Pfaff}, \emph{Jacobi} f\"{u}r
die allgemeine Theorie der Differentialgleichungen so lange
unfruchtbar geblieben sind.  F\"{u}r den gegenw\"{a}rtigen Zweck
gen\"{u}gt es, aus diesem allgemeinen Theile der Lehre von den
ausgedehnten Gr\"{o}ssen, wo weiter nichts vorausgesetzt wird,
als was in dem Begriffe derselben schon enthalten ist, zwei
Punkte hervorzuheben, woven der erste die Erzeugung des Begriffs
einer mehrfach ausgedehnten Mannigfaltigkeit, der zweite die
Zur\"{u}ckf\"{u}hrung der Ortsbestimmungen in einer gegebenen
Mannigfaltigkeit auf Quantit\"{a}tsbestimmungen betrifft und das
wesentliche Kennzeichen einer $n \,$\-fachen Ausdehnung deutlich
machen wird.

\medbreak

\centerline{2.}

\nobreak\medskip

Geht man bei einem Begriffe, dessen Bestimmungsweisen eine
stetige Mannigfaltigkeit bilden, von einer Bestimmungsweise auf
eine bestimmte Art zu einer andern \"{u}ber, so bilden die
durchlaufenen Bestimmungsweisen eine einfach ausgedehnte
Mannigfaltigkeit, deren wesentliches Kennzeichen ist, dass in ihr
von einem Punkte nur nach zwei Seiten, vorw\"{a}rts order
r\"{u}ckw\"{a}rts, ein stetiger Fortgang m\"{o}glich ist.  Denkt
man sich nun, dass diese Mannigfaltigkeit wieder in eine andere,
v\"{o}llig verschiedene, \"{u}bergeht, und zwar wieder auf
bestimmte Art, d.~h.\ so, dass jeder Punkt in einen bestimmten
Punkt der andern \"{u}bergeht, so bilden s\"{a}mmtliche so
erhaltene Bestimmungsweisen eine zweifach ausgedehnte
Mannigfaltigkeit.  In \"{a}hnlicher Weise erh\"{a}lt man eine
dreifach ausgedehnte Mannigfaltigkeit, wenn man sich vorstellt,
dass eine zweifach ausgedehnte in eine v\"{o}llig verschiedene
auf bestimmte Art \"{u}bergeht, und es ist leicht zu sehen, wie
man diese Construction fortsetzen kann.  Wenn man, anstatt den
Begriff als bestimmbar, seinen Gegenstand als ver\"{a}nderlich
betrachtet, so kann diese Construction bezeichnet werden als eine
Zusammensetzung einer Ver\"{a}nderlichkeit von $n + 1$
Dimensionen aus einer Ver\"{a}nderlichkeit von $n$ Dimensionen
und aus einer Ver\"{a}nderlichkeit von Einer Dimension.

\medbreak

\centerline{3.}

\nobreak\medskip

Ich werde nun zeigen, wie man umgekehrt eine
Ver\"{a}nderlichkeit, deren Gebiet gegeben ist, in eine
Ver\"{a}nderlichkeit von einer Dimension und eine
Ver\"{a}nderlichkeit von weniger Dimensionen zerlegen kann.  Zu
diesem Ende denke man sich ein ver\"{a}nderliches St\"{u}ck einer
Mannigfaltigkeit von Einer Dimension---von einem festen
Anfangspunkte an gerechnet, so dass die Werthe desselben unter
einander vergleichbar sind---, welches f\"{u}r jeden Punkt der
gegebenen Mannigfaltigkeit einen bestimmten mit ihm stetig sich
\"{a}ndernden Werth hat, oder mit andern Worten, man nehme
innerhalb der gegebenen Mannigfaltigkeit eine stetige Function
des Orts an, un zwar eine solche Function, welche nicht l\"{a}ngs
eines Theils dieser Mannigfaltigkeit constant ist.  Jedes system
von Punkten, wo die Function einen constanten Werth hat, bildet
dann eine stetige Mannigfaltigkeit von weniger Dimensionen, als
die gegebene.  Diese Mannigfaltigkeiten gehen bei Aenderung der
Function stetig in einander \"{u}ber; man wird daher annehmen
k\"{o}nnen, dass aus einer von ihnen die \"{u}brigen hervorgehen,
und es wird dies, allgemein zu reden, so geschehen k\"{o}nnen,
dass jeder Punkt in einen bestimmten Punkt der andern
\"{u}bergeht; die Ausnahmsf\"{a}lle, deren Untersuchung wichtig
ist, k\"{o}nnen hier unber\"{u}cksichtigt bleiben.  Hierdurch
wird die Ortsbestimmung in der gegebenen Mannigfaltigkeit
zur\"{u}ckgef\"{u}hrt auf eine Gr\"{o}ssenbestimmung und auf eine
Ortsbestimmung in einer minderfach ausgedehnten Mannigfaltigkeit.
Es ist nun leicht zu zeigen, dass diese Mannigfaltigkeit $n - 1$
Dimensionen hat, wenn die gegebene Mannigfaltigkeit eine $n \,$\-fach
ausgedehnte ist.  Durch $n$\-malige Wiederholung dieses
Verfahrens wird daher die Ortsbestimmung in eine $n \,$\-fach
ausgedehnten Mannigfaltigkeit auf $n$ Gr\"{o}ssenbestimmungen,
und also die Ortsbestimmung in einer gegebenen Mannigfaltigkeit,
wenn dieses m\"{o}glich ist, auf eine endliche Anzahl von
Quantit\"{a}tsbestimmungen zur\"{u}ckgef\"{u}hrt.  Es giebt
indess auch Mannigfaltigkeiten, in welchen die Ortsbestimmung
nicht eine endliche Zahl, sondern entweder eine unendliche Reihe
oder eine stetige Mannigfaltigkeit von Gr\"{o}ssenbestimmungen
erfordert.  Solche Mannigfaltigkeiten bilden z.~B. die
m\"{o}glichen Bestimmungen einer Function f\"{u}r ein gegebenes
Gebiet, die m\"{o}glichen Gestalten ein err\"{a}umlichen Figur
u.~s.~w.

\begin{center}
\large\bfseries
II.  Massverh\"{a}ltnisse, deren eine Mannigfaltigkeit von
$n$~Dimensionen f\"{a}hig ist, unter der Voraussetzung, dass die
Linien unabh\"{a}ngig von der Lage eine L\"{a}nge besitzen, also jede
Linie durch jede messbar ist.
\end{center}

Es folgt nun, nachdem der Begriff einer $n \,$\-fach ausgedehnten
Mannigfaltigkeit construirt und als wesentliches Kennzeichen
derselben gefunden worden ist, dass sich die Ortsbestimmung in
derselben auf $n$ Gr\"{o}ssenbestimmungen zur\"{u}ckf\"{u}hren
l\"{a}sst, als zweite der oben gestellten Aufgaben eine
Untersuchung \"{u}ber die Massverh\"{a}ltnisse, deren ein solche
Mannigfaltigkeit f\"{a}hig ist, und \"{u}ber die Bedingungen,
welche zur Bestimmung dieser Massverh\"{a}ltnisse hinreichen.
Diese Massverh\"{a}ltnisse lassen sich nur in abstracten
Gr\"{o}ssenbegriffen untersuchen und im Zusammenhange nur durch
Formeln darstellen; unter gewissen Voraussetzungen kann man sie
indess in Verh\"{a}ltnisse zerlegen, welche einzeln genommen
einer geometrischen Darstellung f\"{a}hig sind, und hiedurch wird
es m\"{o}glich, die Resultate der Rechnung geometrisch
auszudr\"{u}cken.  Es wird daher, um festen Boden zu gewinnen,
zwar eine abstracte Untersuchung in Formeln nicht zu vermeiden
sein, die Resultate derselben aber werden sich im geometrischen
Gewande darstellen lassen.  Zu Beidem sind die Grundlagen
enthalten in der ber\"{u}hmten Abhandlung des Herrn Geheimen
Hofraths \emph{Gauss} \"{u}ber die krummen Fl\"{a}chen.

\medbreak

\centerline{1.}

\nobreak\medskip

Massbestimmungen erfordern eine Unabh\"{a}ngigkeit der
Gr\"{o}ssen vom Ort, die in mehr als einer Weise stattfinden
kann; die zun\"{a}chst sich darbietende Annahme, welche ich hier
verfolgen will, ist wohl die, dass die L\"{a}nge der Linien
unabh\"{a}ngig von der Lage sei, also jede Linie durch jede
messbar sei.  Wird die Ortsbestimmung auf Gr\"{o}ssenbestimmungen
zur\"{u}ckgef\"{u}hrt, also die Lage eines Punktes in der
gegebenen $n \,$\-fach ausgedehnten Mannigfaltigkeit durch $n$
ver\"{a}nderliche Gr\"{o}ssen $x_1, x_2, x_3$, und so fort bis
$x_n$ ausgedr\"{u}ckt, so wird die Bestimmung einer Linie darauf
hinauskommen, dass die Gr\"{o}ssen~$x$ als Functionen Einer
Ver\"{a}nderlichen gegeben werden.  Die Aufgabe ist dann, f\"{u}r
die L\"{a}nge der Linien einen mathematischen Ausdruck
aufzustellen, zu welchem Zwecke die Gr\"{o}ssen $x$ als in
Einheiten ausdr\"{u}ckbar betrachtet werden m\"{u}ssen.  Ich
werde diese Aufgabe nur unter gewissen Beschr\"{a}nkungen
behandeln und beschr\"{a}nke mich erstlich auf solche Linien, in
welchen die Verh\"{a}ltnisse zwischen den Gr\"{o}ssen $dx$---den
zusammengeh\"{o}rigen Aenderungen der Gr\"{o}ssen $x$---sich
stetig \"{a}ndern; man kann dann die Linien in Elemente zerlegt
denken, innerhalb deren die Verh\"{a}ltnisse der Gr\"{o}ssen $dx$
als constant betrachtet werden d\"{u}rfen, und die Aufgabe kommt
dann darauf zur\"{u}ck, f\"{u}r jeden Punkt einen allgemeinen
Ausdruck des von ihm ausgeh\-enden Linienelements $ds$
aufzustellen, welcher also die Gr\"{o}ssen $x$ und die
Gr\"{o}ssen $dx$ enthalten wird.  Ich nehme nun zweitens an, dass
die L\"{a}nge des Linienelements, von Gr\"{o}ssen zweiter Ordnung
abgesehen, unge\"{a}ndert bleibt, wenn s\"{a}mmtliche Punkte
desselben dieselbe unendlich kleine Orts\"{a}nderung erleiden,
worin zugleich enthalten ist, dass, wenn s\"{a}mmtliche
Gr\"{o}ssen $dx$ in demselben Verh\"{a}ltnisse wachsen, das
Linienelement sich ebenfalls in diesem Verh\"{a}ltnisse
\"{a}ndert.  Unter diesen Annahmen wird das Linienelement eine
beliebige homogene Function ersten Grades der Gr\"{o}ssen $dx$
sein k\"{o}nnen, welche unge\"{a}ndert bleibt, wenn
s\"{a}mmtliche Gr\"{o}ssen $dx$ ihr Zeichen \"{a}ndern, und worin
die willk\"{u}rlichen Constanten stetige Functionen der
Gr\"{o}ssen $x$ sind.  Um die einfachsten F\"{a}lle zu finden,
suche ich zun\"{a}chst einen Ausdruck f\"{u}r die ${(n-1)}$\-fach
ausgedehnten Mannigfaltigkeiten, welche vom Anfangspunkte des
Linienelements \"{u}berall gleich weit abstehen, d.~h.\ ich suche
eine stetige Function des Orts, welche sie von einander
unterscheidet.  Diese wird vom Anfangspunkt aus nach allen Seiten
entweder ab- oder zunehmen m\"{u}ssen; ich will annehmen, dass
sie nach allen Seiten zunimmt und also in dem Punkte ein Minimum
hat.  Es muss dann, wenn ihre ersten und zweiten
Differentialquotienten endlich sind, das Differential erster
Ordnung verschwinden und das zweiter Ordnung darf nie negativ
werden; ich nehme an, dass es immer positiv bleibt.  Dieser
Differentialausdruck zweiter Ordnung bleibt alsdann constant,
wenn $ds$ constant bleibt, und w\"{a}chst im quadratischen
Verh\"{a}ltnisse, wenn die Gr\"{o}ssen $dx$ und also auch $ds$
sich s\"{a}mmtlich in demselben Verh\"{a}ltnisse \"{a}ndern; er
ist also $=$ const.~$ds^2$ und folglich ist $ds =$ der
Quadratwurzel aus einer immer positiven ganzen homogenen Function
zweiten Grades der Gr\"{o}ssen $dx$, in welcher die Coefficienten
stetige Functionen der Gr\"{o}ssen $x$ sind.  F\"{u}r den Raum
wird, wenn man die Lage der Punkte durch rechtwinklige
Coordinaten ausdr\"{u}ckt,
$ds = \sqrt{ \sum (dx)^2 }$;
der Raum ist also unter diesem einfachsten Falle enthalten.  Der
n\"{a}chst einfache Fall w\"{u}rde wohl die Mannigfaltigkeiten
umfassen, in welchen sich das Linienelement durch die vierte
Wurzel aus einem Differentialausdrucke vierten Grades
ausdr\"{u}cken l\"{a}sst.  Die Untersuchung dieser allgemeinern
Gattung w\"{u}rde zwar keine wesentlich andere Principien
erfordern, aber ziemlich zeitraubend sein und
verh\"{a}ltnissm\"{a}ssig auf die Lehre vom Raume wenig neues
Licht werfen, zumal da sich die Resultate nicht geometrisch
ausdr\"{u}cken lassen; ich beschr\"{a}nke mich daher auf die
Mannigfaltigkeiten, wo das Linienelement durch die Quadratwurzel
aus einem Differentialausdruck zweiten Grades ausgedr\"{u}ckt
wird.  Man kann einen solchen Ausdruck in einen andern
\"{a}hnlichen transformiren, indem man f\"{u}r die $n$
unabh\"{a}ngigen Ver\"{a}nderlichen Functionen von $n$ neuen
unabh\"{a}ngigen Ver\"{a}nderlichen setzt.  Auf diesem Wege wird
man aber nicht jeden Ausdruck in jeden transformiren k\"{o}nnen;
den der Ausdruck enth\"{a}lt
$\displaystyle n \frac{n+1}{2}$
Coefficienten, welche willk\"{u}rliche Functionen der
unabh\"{a}ngigen Ver\"{a}nderlichen sind; durch Einf\"{u}hrung
neuer Ver\"{a}nderlicher wird man aber nur $n$ Relationen
gen\"{u}gen und also nur $n$ der Coefficienten gegebenen
Gr\"{o}ssen gleich machen k\"{o}nnen.  Es sind dann die
\"{u}brigen
$\displaystyle n \frac{n-1}{2}$
durch die Natur der darzustellenden Mannigfaltigkeit schon
v\"{o}llig bestimmt, und zur Bestimmung ihrer
Massverh\"{a}ltnisse also $\displaystyle n \frac{n-1}{2}$
Functionen des Orts erforderlich.  Die Mannigfaltigkeiten, in
welchen sich, wie in der Ebene und im Raume, das Linienelement
auf die Form
$\sqrt{ \sum dx^2 }$
bringen l\"{a}sst, bilden daher nur einen besondern Fall der hier
zu untersuchenden Mannigfaltigkeiten; sie verdienen wohl einen
besonderen Namen, und ich will also diese Mannigfaltigkeiten, in
welchen sich das Quadrat des Linienelements auf die Summe der
Quadrate von selbst\"{a}ndigen Differentialien bringen l\"{a}sst,
eben nennen.  Um nun die wesentlichen Verschiedenheiten
s\"{a}mmtlicher in der vorausgesetzten Form darstellbarer
Mannigfaltigkeiten \"{u}bersehen zu k\"{o}nnen, ist es
n\"{o}thig, die von der Darstellungsweise herr\"{u}hrenden zu
beseitigen, was durch Wahl der ver\"{a}nderlichen Gr\"{o}ssen
nach einem bestimmten Princip erreicht wird.

\medbreak

\centerline{2.}

\nobreak\medskip

Zu diesem Ende denke man sich von einem beliebigen Punkte aus das
System der von ihm ausgehenden k\"{u}rzesten Linien constuirt;
die Lage eines unbestimmten Punktes wird dann bestimmt werden
k\"{o}nnen durch die Anfangsrichtung der k\"{u}rzesten Linie, in
welcher er liegt, und durch seine Entfernung in derselben vom
Anfangspunke und kann daher durch die Verh\"{a}ltnisse der
Gr\"{o}ssen $dx^0$, d.~h.\ der Gr\"{o}ssen $dx$ in Anfang dieser
k\"{u}rzesten Linie und durch die L\"{a}nge~$s$ dieser Linie
ausgedr\"{u}ckt werden.  Man f\"{u}hre nun statt $dx^0$ solche
aus ihnen gebildete lineare Ausdr\"{u}cke $d\alpha$ ein, dass der
Anfangswerth des Quadrats des Linienelements gleich der Summe der
Quadrate dieser Ausdr\"{u}cke wird, so dass die unabh\"{a}ngigen
Variabeln sind: die Gr\"{o}sse $s$ und die Verh\"{a}ltnisse der
Gr\"{o}ssen $d\alpha$; und setze schliesslich statt d$\alpha$
solche ihnen proportionale Gr\"{o}ssen $x_1, x_2,\ldots, x_n$,
dass die Quadratsumme $= s^2$ wird.  F\"{u}hrt man diese
Gr\"{o}ssen ein, so wird f\"{u}r unendlich kleine Werthe von $x$
das Quadrat des Linienelement $= \sum dx^2$, das Glied der
n\"{a}chsten Ordnung in demselben aber gleich einem homogenen
Ausdruck zweiten Grades der
$\displaystyle n \frac{n-1}{2}$
Gr\"{o}ssen
$(x_1 \, dx_2 - x_2 \, dx_1)$,
$(x_1 \, dx_3 - x_3 \,dx_1),\ldots,$
also eine unendliche kleine Gr\"{o}sse von der vierten Dimension,
so dass man eine endliche Gr\"{o}sse erh\"{a}lt, wenn man sie
durch das Quadrat des unendlich kleinen Dreiecks dividirt, in
dessen Eckpunkten die Werthe der Ver\"{a}nderlichen sind
$(0,0,0,\ldots)$,
$(x_1, x_2, x_3,\ldots)$,
$(dx_1, dx_2, dx_3,\ldots)$.
Diese Gr\"{o}sse beh\"{a}lt denselben Werth, so lange die
Gr\"{o}ssen $x$ und $dx$ in denselben bin\"{a}ren Linearformen
enthalten sind, oder so lange die beiden k\"{u}rzesten Linien von
den Werthen $0$ bis zu den Werthen~$x$ und von den Werthen $0$
bis zu den Werthen $dx$ in demselben Fl\"{a}chenelement bleiben,
und h\"{a}ngt also nur von Ort und Richtung desselben ab.  Sie
wird offenbar $= 0$, wenn die dargestellte Mannigfaltigkeit eben,
d.~h.\ das Quadrat des Linienelements auf $\sum dx^2$ reducirbar
ist, und kann daher als das Mass der in diesem Punkte in dieser
Fl\"{a}chenrichtung stattfindenden Abweichung der
Mannigfaltigkeit von der Ebenheit angesehen werden.  Multiplicirt
mit $-\frac{3}{4}$ wird sie der Gr\"{o}sse gleich, welche Herr
Geheimer Hofrath \emph{Gauss} das Kr\"{u}mmungsmass einer
Fl\"{a}che ganannt hat.  Zur Bestimmung der Massverh\"{a}ltnisse
einer $n \,$\-fach ausgedehnten in der vorausgesetzten Form
darstellbaren Mannigfaltigkeit wurden vorhin
$\displaystyle n \frac{n-1}{2}$
Functionen des Orts n\"{o}thig gefunden; wenn also das
Kr\"{u}mmungsmass in jedem Punkte in
$\displaystyle n \frac{n-1}{2}$
Fl\"{a}chenrichtungen gegeben wird, so werden daraus die
Massverh\"{a}ltnisse der Mannigfaltigkeit sich bestimmen lassen,
wofern nur zwischen diesen Werthen keine identischen Relationen
stattfinden, was in der That, allgemein zu reden, nicht der Fall
ist.  Die Massverh\"{a}ltnisse dieser Mannigfaltigkeiten, wo das
Linienelement durch die Quadratwurzel aus einem
Differentialausdruck zweiten Grades dargestellt wird, lassen sich
so auf eine von der Wahl der ver\"{a}nderlichen Gr\"{o}ssen
v\"{o}llig unabh\"{a}ngige Weise ausdr\"{u}cken.  Ein ganz
\"{a}hnlicher Weg l\"{a}sst sich zu diesem Ziele auch bei den
Mannigfaltigkeiten einschlagen, in welchen das Linienelement
durch einen weniger einfachen Ausdruck, z.~B. durch die vierte
Wurzel aus einem Differentialausdruck vierten Grades,
ausgedr\"{u}ckt wird.  Es w\"{u}rde sich dann das Linienelement,
allgemein zu reden, nicht mehr auf die Form der Quadratwurzel aus
einer Quadratsumme von Differentialausdr\"{u}cken bringen lassen
und also in dem Ausdrucke f\"{u}r das Quadrat des Linienelements
die Abweichung von der Ebenheit eine unendlich kleine Gr\"{o}sse
von der zweiten Dimension sein, w\"{a}hrend sie bei jenen
Mannigfaltigkeiten eine unendlich kleine Gr\"{o}sse von der
vierten Dimension war.  Diese Eigenth\"{u}mlichkeit der letztern
Mannigfaltigkeiten kann daher wohl Ebenheit in den kleinsten
Theilen genannt werden.  Die f\"{u}r den jetzigen Zweck
wichtigste Eigenth\"{u}mlichkeit dieser Mannigfaltigkeiten,
derentwegen sie hier allein untersucht worden sind, ist aber die,
dass sich die Verh\"{a}ltnisse der zweifach ausgedehnten
geometrisch durch Fl\"{a}chen darstellen und die der mehrfach
ausgedehnten auf die der in ihnen enthaltenen Fl\"{a}chen
zur\"{u}ckf\"{u}hren lassen, was jetzt noch einer kurzen
Er\"{o}rterung bedarf.

\medbreak

\centerline{3.}

\nobreak\medskip

In die Auffassung der Fl\"{a}chen mischt sich neben den inneren
Massverh\"{a}ltnissen, bei welchen nur die L\"{a}nge der Wege in
ihnen in Betracht kommt, immer auch ihre Lage zu ausser ihnen
gelegenen Punkten.  Man kann aber von den \"{a}ussern
Verh\"{a}ltnissen abstrahiren, indem man solche Ver\"{a}nderungen
mit ihnen vornimmt, bei denen die L\"{a}nge der Linien in ihnen
unge\"{a}ndert bleibt, d.~h.\ sie sich beliebig---ohne
Dehnung---gebogen denkt, und alle so auseinander entstehenden
Fl\"{a}chen als gleichartig betrachtet.  Es gelten also
z.~B.\ beliebige cylindrische oder conische Fl\"{a}chen einer
Ebene gleich, weil sie sich durch blosse Biegung aus ihr bilden
lassen, wobei die innern Massverh\"{a}ltnisse bleiben, und
s\"{a}mmtliche S\"{a}tze \"{u}ber dieselben---also die ganze
Planimetrie---ihre G\"{u}ltigkeit behalten; dagegen gelten sie
als wesentlich verschieden von der Kugel, welche sich nicht ohne
Dehnung in eine Ebene verwandeln l\"{a}sst.  Nach dem vorigen
Untersuchung werden in jedem Punkte die innern
Massverh\"{a}ltnisse einer zweifach ausgedehnten Gr\"{o}sse, wenn
sich das Linienelement durch die Quadratwurzel aus einem
Differentialausdruck zweiten Grades ausdr\"{u}cken l\"{a}sst, wie
dies bei den Fl\"{a}chen der Fall ist, charakterisirt durch das
Kr\"{u}mmungsmass.  Dieser Gr\"{o}sse l\"{a}sst sich nun bei den
Fl\"{a}chen die anschauliche Bedeutung geben, dass sie das
Product aus den beiden Kr\"{u}mmungen der Fl\"{a}che in diesem
Punkte ist, oder auch, dass das Product derselben in ein
unendlich kleines aus k\"{u}rzesten Linien gebildetes Dreieck
gleich ist dem halben Ueberschusse seiner Winkelsumme \"{u}ber
zwei Rechte in Theilen des Halbmessers.  Die erste Definition
w\"{u}rde den Satz voraussetzen, dass das Product der beiden
Kr\"{u}mmungshalbmesser bei der blossen Biegung einer Fl\"{a}che
unge\"{a}ndert bleibt, die zweite, dass an demselben Orte der
Ueberschuss der Winkelsumme eines unendlich kleinen Dreiecks
\"{u}ber zwei Rechte seinem Inhalte proportional ist.  Um die
Kr\"{u}mmungsmass einer $n \,$\-fach ausgedehnten Mannigfaltigkeit in
einem gegebenen Punkte und einer gegebenen durch ihn gelegten
Fl\"{a}chenrichtung eine greifbare Bedeutung zu geben, muss man
davon ausgehen, dass eine von einem Punkte ausgehende
k\"{u}rzeste Linie v\"{o}llig bestimmt ist, wenn ihre
Ansfangsrichtung gegeben ist.  Hienach wird man eine bestimmte
Fl\"{a}che erhalten, wenn man s\"{a}mmtliche von dem gegebenen
Punkte ausgehenden und in dem gegebenene Fl\"{a}chenelement
liegenden Anfangsrichtungen zu k\"{u}rzesten Linien
verl\"{a}ngert, und diese Fl\"{a}che hat in dem gegebenen Punkte
ein bestimmtes Kr\"{u}mmungsmass, welches zugleich das
Kr\"{u}mmungsmass der $n \,$\-fach ausgedehnten Mannigfaltigkeit in
dem gegebenen Punkte und der gegebenen Fl\"{a}chenrichtung ist.

\medbreak

\centerline{4.}

\nobreak\medskip

Es sind nun noch, ehe die Anwendung auf den Raum gemacht wird,
einige Betrachtungen \"{u}ber die ebenen Mannigfaltigkeiten im
Allgemeinen n\"{o}thig, d.~h.\ \"{u}ber diejenigen, in welchen
das Quadrat des Linienelements durch eine Quadratsumme
vollst\"{a}ndiger Differentialien darstellbar ist.

In einer ebenen $n \,$\-fach ausgedehnten Mannigfaltigkeit ist das
Kr\"{u}mmungsmass in jedem Punkte in jeder Richtung Null; es
reicht aber nach der fr\"{u}hern Untersuchung, um die
Massverh\"{a}ltnisse zu bestimmen, hin zu wissen, dass es in
jedem Punkte in $\displaystyle n \frac{n-1}{2}$
Fl\"{a}chenrichtungen, deren Kr\"{u}mmungsmasse von einander
unabh\"{a}ngig sind, Null sei.  Die Mannigfaltigkeiten, deren
Kr\"{u}mmungsmass \"{u}berall $= 0$ ist, lassen sich betrachten
als ein besonderer Fall derjenigen Mannigfaltigkeiten, deren
Kr\"{u}mmungsmass allenthalben constant ist.  Der gemeinsame
Charakter dieser Mannigfaltigkeiten, deren Kr\"{u}mmungsmass
constant ist, kann auch so ausgedr\"{u}ckt werden, dass sich die
Figuren in ihnen ohne Dehnung bewegen lassen.  Denn offenbar
w\"{u}rden die Figuren in ihnen nicht beliebig verschiebbar und
drehbar sein k\"{o}nnen, wenn nicht in jedem Punkte in allen
Richtungen das Kr\"{u}mmungsmass dasselbe w\"{a}re.  Andererseits
aber sind durch das Kr\"{u}mmungsmass die Massverh\"{a}ltnisse
der Mannigfaltigkeit vollst\"{a}ndig bestimmt; es sind daher um
einen Punkt nach allen Richtungen die Massverh\"{a}ltnisse genau
dieselben, wie um einen andern, und also von ihm aus dieselben
Constructionen ausf\"{u}hrbar, und folglich kann in den
Mannigfaltigkeiten mit constantem Kr\"{u}mmungsmass den Figuren
jede beliebige Lage gegeben werden.  Die Massverh\"{a}ltnisse
dieser Mannigfaltigkeiten h\"{a}ngen nur von dem Werthe des
Kr\"{u}mmungsmasses ab, und in Bezung auf die analytishe
Darstellung mag bemerkt werden, dass, wenn man diesen Werth durch
$\alpha$ bezeichnet, dem Ausdruck f\"{u}r das Linienelement die
Form
\[ \frac{1}{\displaystyle 1 + \frac{\alpha}{4} {\textstyle\sum} x^2}
   \sqrt{ {\textstyle\sum} dx^2 } \]
gegeben werden kann.

\medbreak

\centerline{5.}

\nobreak\medskip

Zur geometrischen Erl\"{a}uterung kann die Betrachtung der
\emph{Fl\"{a}chen} mit constantem Kr\"{u}mmungsmass dienen.  Es
ist leicht zu sehen, dass sich die Fl\"{a}chen, deren
Kr\"{u}mmungsmass positiv ist, immer auf eine Kugel, deren Radius
gleich~$1$ dividirt durch die Wurzel aus dem Kr\"{u}mmungsmass
ist, wickeln lassen werden; um aber die ganze Mannigfaltigkeit
dieser Fl\"{a}chen zu \"{u}bersehen, gebe man einer derselben die
Gestalt einer Kugel und den \"{u}brigen die Gestalt von
Umdrehungsfl\"{a}chen, welche sie im Aequator ber\"{u}hren.  Die
Fl\"{a}chen mit gr\"{o}sserem Kr\"{u}mmungsmass, als diese Kugel,
werden dann die Kugel von innen ber\"{u}hren und eine Gestalt
annehmen, wie der \"{a}ussere der Axe abgewandte Theil der
Oberfl\"{a}che eines Ringes; sie w\"{u}rden sich auf Zonen von
Kugeln mit kleinerem Halbmesser wickeln lassen, aber mehr als
einmal herumreichen.  Die Fl\"{a}chen mit kleinerem positiven
Kr\"{u}mmungsmass wird man erhalten, wenn man aus Kugelfl\"{a}chen
mit gr\"{o}sserem Radius ein von zwei gr\"{o}ssten Halbkreisen
begrenztes St\"{u}ck ausschneidet und die Schnittlinien
zusammenf\"{u}gt.  Die Fl\"{a}che mit dem Kr\"{u}mmungsmass Null
wird eine auf dem Aequator stehende Cylinderfl\"{a}che sein; die
Fl\"{a}chen mit negativem Kr\"{u}mmungsmass aber werden diesen
Cylinder von aussen ber\"{u}hren und wie der innere der Axe
zugewandte Theil der Oberfl\"{a}che eines Ringes geformt sein.
Denkt man sich diese Fl\"{a}che als Ort f\"{u}r in ihnen
bewegliche Fl\"{a}chenst\"{u}cke, wie den Raum als Ort f\"{u}r
K\"{o}rper, so sind in allen diesen Fl\"{a}chen die
Fl\"{a}chenst\"{u}cke ohne Dehnung beweglich.  Die Fl\"{a}chen
mit positivem Kr\"{u}mmungsmass lassen sich stets so formen, dass
die Fl\"{a}chenst\"{u}cke auch ohne Biegung beliebig bewegt
werden k\"{o}nnen, n\"{a}mlich zu Kugelfl\"{a}chen, die mit
negativem aber nicht.  Ausser dieser Unabh\"{a}ngigkeit der
Fl\"{a}chenst\"{u}cke vom Ort findet bei der Fl\"{a}che mit dem
Kr\"{u}mmungsmass Null auch eine Unabh\"{a}ngigkeit der Richtung
vom Ort statt, welche bei den \"{u}brigen Fl\"{a}chen nicht
stattfindet.

\begin{center}
\large\bfseries
III.  Anwendung auf den Raum.
\end{center}

\centerline{1.}

\nobreak\medskip

Nach diesen Untersuchungen \"{u}ber die Bestimmung der
Massverh\"{a}ltnisse einer $n \,$\-fach ausgedehnten Gr\"{o}sse lassen
sich nun die Bedingungen angeben, welche zur Bestimmung der
Massverh\"{a}ltnisse des Raumes hinreichend und nothwendig sind,
wenn Unabh\"{a}ngigkeit der Linien von der Lage und
Darstellbarkeit des Linienelements durch die Quadratwurzel aus
einem Differentialausdrucke zweiten Grades, also Ebenheit in den
kleinsten Theilen vorausgesetzt wird.

Sie lassen sich erstens so ausdr\"{u}cken, dass das
Kr\"{u}mmungsmass in jedem Punkte in drei Fl\"{a}chenrichtungen
$= 0$ ist, und es sind daher die Massverh\"{a}ltnisse des Raumes
bestimmt, wenn die Winkelsumme im Dreieck allenthalben gleich
zwei Rechten ist.

Setzt man aber zweitens, wie \emph{Euklid}, nicht bloss eine
von der Lage unabh\"{a}ngige Existenz der Linien, sondern auch
der K\"{o}rper voraus, so folgt, dass das Kr\"{u}mmungsmass
allenthalben constant ist, und es ist dann in allen Dreiecken die
Winkelsumme bestimmt, wenn sie in Einem bestimmt ist.

Endlich k\"{o}nnte man drittens, anstatt die L\"{a}nge der Linien
als unabh\"{a}ngig von Ort und Richtung anzunehmen, auch eine
Unabh\"{a}ngigkeit ihre L\"{a}nge und Richtung vom Ort
voraussetzen.  Nach dieser Auffassung sind die Orts\"{a}nderungen
oder Ortsverschiedenheiten complexe in drei unabh\"{a}ngige
Einheiten ausdr\"{u}ckbare Gr\"{o}ssen.

\medbreak

\centerline{2.}

\nobreak\medskip

Im Laufe der bisherigen Betrachtungen wurden zun\"{a}chst die
Aus\-dehn\-ungs- oder Gebietsverh\"{a}ltnisse von den
Massverh\"{a}ltnissen gesondert, und gefunden, dass bei denselben
Ausdehnungsverh\"{a}ltnissen verschiedene Massverh\"{a}ltnisse
denkbar sind; es wurden dann die Systeme einfacher
Massbestimmungen aufgesucht, durch welche die
Massverh\"{a}ltnisse des Raumes v\"{o}llig bestimmt sind und von
welchen alle S\"{a}tze \"{u}ber dieselben eine nothwendige Folge
sind; es bleibt nun die Frage zu er\"{o}rtern, wie, in welchem
Grade und in welchem Unfange diese Voraussetzungen durch die
Erfahrung verb\"{u}rgt werden.  In dieser Beziehung findet
zwischen den blossen Ausdehnungsverh\"{a}ltnissen und den
Massverh\"{a}ltnissen eine wesentliche Verschiedenheit statt,
insofern bei erstern, wo die m\"{o}glichen F\"{a}lle eine
discrete Mannigfaltigkeit bilden, die Aussagen der Erfahrung zwar
nie v\"{o}llig gewiss, aber nicht ungenau sind, w\"{a}hrend bei
letztern, wo die m\"{o}glichen F\"{a}lle eine stetige
Mannigfaltigkeit bilden, jede Bestimmung aus der Erfahrung immer
ungenau bleibt---es mag die Wahrscheinlichkeit, dass si nahe
richtig ist, noch so gross sein.  Dieser Umstand wird wichtig bei
der Ausdehnung dieser empirischen Bestimmungen \"{u}ber die
Grenzen der Beobachtung in's Unmessbargrosse und Unmessbarkleine;
denn die letztern k\"{o}nnen offenbar jenseits der Grenzen der
Beobachtung immer ungenauer werden, die ersteren aber nicht.

Bei der Ausdehnung der Raumconstructionen in's Unmessbargrosse
ist Unbegrenztheit und Unendlichkeit zu scheiden; jene geh\"{o}rt
zu den Ausdehnungsverh\"{a}ltnissen, diese zu den
Massverh\"{a}ltnissen.  Dass der Raum eine unbegrenzte dreifach
ausgedehnte Mannigfaltigkeit sei, ist eine Voraussetzung, welche
bei jeder Auffassung der Aussenwelt angewandt wird, nach welcher
in jedem Augenblicke das Gebiet der wirklichen Wahrnehmungen
erg\"{a}nzt und die m\"{o}glichen Orte eines gesuchten
Gegenstandes construirt werden und welche sich bei diesen
Anwendungen fortw\"{a}hrend best\"{a}tigt.  Die Unbegrenztheit
des Raumes besitzt daher eine gr\"{o}ssere empirische Gewissheit,
als irgend eine \"{a}ussere Erfahrung.  Hieraus folt aber die
Unendlichkeit keineswegs; vielmehr w\"{u}rde der Raum, wenn man
Unabh\"{a}ngigkeit der K\"{o}rper vom Ort voraussetzt, ihm also
ein constantes Kr\"{u}mmungsmass zuschreibt, nothwendig endlich
sein, so bald dieses Kr\"{u}mmungsmass einen noch so kleinen
positiven Werth h\"{a}tte.  Man w\"{u}rde, wenn man die in einem
Fl\"{a}chenelement liegenden Anfangsrichtungen zu k\"{u}rzesten
Linien verl\"{a}ngert, eine unbegrenzte Fl\"{a}che mit constantem
positiven Kr\"{u}mmungsmass, also eine Fl\"{a}che erhalten,
welche in einer ebenen dreifach ausgedehnten Mannigfaltigkeit die
Gestalt einer Kugelfl\"{a}che annehmen w\"{u}rde und welche
folglich endlich ist.

\medbreak

\centerline{3.}

\nobreak\medskip

Die Fragen \"{u}ber das Unmessbargrosse sind f\"{u}r die
Naturerkl\"{a}rung m\"{u}ssige Fragen.  Anders verh\"{a}lt es
sich aber mit den Fragen \"{u}ber das unmessbarkleine.  Auf der
Genauigkeit, mit welcher wir die Erscheinungen in's
Unendlichkleine verfolgen, beruht wesentlich die Erkenntniss
ihres Causalzusammenhangs.  Die Fortschritte der letzten
Jahrhunderte in der Erkenntniss der mechanischen Natur sind fast
allein bedingt durch die Genauigkeit der Construction, welche
durch die Erfindung der Analysis des Unendlichen und die von
\emph{Archimed}, \emph{Galil\"{a}i} und \emph{Newton}
aufgefundenen einfachen Grundbegriffe, deren sich die heutige
Physik bedient, m\"{o}glich geworden ist.  In den
Naturwissenschaften aber, wo die einfachen Grundbegriffe zu
solchen Constructionen bis jetzt fehlen, verfolgt man, um den
Causalzusammenhang zu erkennen, die Erscheinungen in's
r\"{a}umlich Kleine, so weit es das Mikroskop nur gestattet.  Die
Fragen \"{u}ber die Massverh\"{a}ltnisse des Raumes im
Unmessbarkleinen geh\"{o}ren also nicht zu den m\"{u}ssigen.

Setzt man voraus, dass die K\"{o}rper unabh\"{a}ngig vom Ort
existiren, so ist dass Kr\"{u}mmungsmass \"{u}berall constant,
und es folgt dann aus den astronomischen Messungen, dass es nicht
von null verschieden sein kann; jedenfalls m\"{u}sste sein
reciprocer Werth eine Fl\"{a}che sein, gegen welche das unsern
Teleskopen zug\"{a}ngliche Gebiet verschwinden m\"{u}sste.  Wenn
aber eine solche Unabh\"{a}ngigkeit der K\"{o}rper vom Ort nicht
stattfindet, so kann man aus den Massverh\"{a}ltnissen im Grossen
nicht auf die im Unendlichkleinen schliessen; es kann dann in
jedem Punkte das Kr\"{u}mmungsmass in drei Richtungen einen
beliebigen Werth haben, wenn nur die ganze Kr\"{u}mmung jedes
messbaren Raumtheils nicht merklich von Null verschieden ist; noch
complicirtere Verh\"{a}ltnisse k\"{o}nnen eintreten, wenn die
vorausgesetzte Darstellbarkeit eines Linienelements durch die
Quadratwurzel aus einem Differentialausdruck zweiten Grades nicht
stattfindet.  Nun scheinen aber die empirischen Begriffe, in
welchen die r\"{a}umlichen Massbestimmungen gegr\"{u}ndet sind,
der Begriff des festen K\"{o}rpers und des Lichtstrahls, im
Unendlichkleinen ihre G\"{u}ltigkeit zu verlieren; es ist also
sehr wohl denkbar, dass die Massverh\"{a}ltnisse des Raumes im
Unendlichkleinen den Voraussetzungen der Geometrie nicht
gem\"{a}ss sind, und dies w\"{u}rde man in der That annehmen
m\"{u}ssen, sobald sich dadurch die Erscheinungen auf einfachere
Weise erkl\"{a}ren liessen.

Die Frage \"{u}ber die G\"{u}ltigkeit der Voraussetzungen der
Geometrie im Unendlichkleinen h\"{a}ngt zusammen mit der Frage
nach dem innern Grunde der Massverh\"{a}ltnisse des Raumes.  Bei
dieser Frage, welche wohl noch zur Lehre vom Raume gerechnet
werden darf, kommt die obige Bemerkung zur Anwendung, dass bei
einer discreten Mannigfaltigkeit das Princip der
Massverh\"{a}ltnisse schon in dem Begriffe dieser
Mannigfaltigkeit enthalten ist, bei einer stetigen aber anders
woher hinzukommen muss.   Es muss also entweder das dem Raume zu
Grunde liegende Wirkliche eine discrete Mannigfaltigkeit bilden,
oder der Grund der Massverh\"{a}ltnisse ausserhalb, in darauf
wirkenden bindenen Kr\"{a}ften, gesucht werden.

Die Entscheidung dieser Fragen kann nur gefunden werden, indem
man von der bisherigen durch die Erfahrung bew\"{a}hrten
Auffassung der Erscheinungen, wozu \emph{Newton} den Grund
gelegt, ausgeht und diese durch Thatsachen, die sich aus ihr
nicht erkl\"{a}ren lassen, getrieben allm\"{a}hlich umarbeitet;
solche Untersuchungen, welche, wie die hier gef\"{u}hrte, von
allgemeinen Begriffen ausgehen, k\"{o}nnen nur dazu dienen, dass
diese Arbeit nicht durch die Beschr\"{a}nktheit der Begriffe
gehindert und der Fortschritt im Erkennen des Zusammenhangs der
Dinge nicht durch \"{u}berlieferte Vorurtheile gehemmt wird.

Es f\"{u}hrt dies hin\"{u}ber in das Gebiet einer andern
Wissenschaft, in das Gebiet der Physik, welches wohl die Natur
der heutigen Veranlassung nicht zu betreten erlaubt.

\bigbreak

\begin{center}
\large\bfseries
Uebersicht
\end{center}

\begin{description}
\item[]
Plan der Untersuchung.

\item[\textmd{I.}]
Begriff einer $n$~fach ausgedehnten Gr\"{o}sse.

\begin{description}
\item[\textmd{\S.~1.}]
Stetige und discrete Mannigfaltigkeiten.  Bestimmte Theile einer
Mannigfaltigkeit heissen Quanta.  Eintheilung der Lehre von den
stetigen Gr\"{o}ssen in die Lehre
\begin{description}
\item[\textmd{1)}]
von den blossen Gebietsverh\"{a}ltnissen, bei welcher eine
Unabh\"{a}ngigkeit der Gr\"{o}ssen vom Ort nicht vorausgesetzt
wird,

\item[\textmd{2)}]
von den Massverh\"{a}ltnissen, bei welcher eine solche
Unabh\"{a}ngigkeit vorausgesetzt werden muss.
\end{description}

\item[\textmd{\S.~2.}]
Erzeugung des Begriffs einer einfach, zweifach,$\ldots$, $n$~fach
ausgedehnten Mannigfaltigkeit.

\item[\textmd{\S.~3.}]
Zur\"{u}ckf\"{u}hrung der Ortsbestimmung in einer gegebenen
Mannigfaltigkeit auf Quantit\"{a}tsbestimmungen.  Wesentliches
Kennzeichen einer $n$~fach ausgedehnten Mannigfaltigkeit.
\end{description}

\item[\textmd{II.}]
Massverh\"{a}ltnisse, deren eine Mannigfaltigkeit von
$n$~Dimensionen f\"{a}hig ist, unter der Voraussetzung, dass die
Linien unabh\"{a}ngig von der Lage eine L\"{a}nge besitzen, also
jede Linie durch jede messbar ist.

\begin{description}
\item[\textmd{\S.~1.}]
Ausdruck des Linienelements.  Als eben werden solche
Mannigfaltigkeiten betrachtet, in denen das Linienelement durch
die Wurzel aus einer Quadratsumme vollst\"{a}ndiger
Differentialien ausdr\"{u}ckbar ist.

\item[\textmd{\S.~2.}]
Untersuchung der $n$~fach ausgedehnten Mannigfaltigkeiten, in
welchen das Linienelement durch die Quadratwurzel aus einem
Differentialausdruck zweiten Grades dargestellt werden kann.  Mass
ihrer Abweichung von der Ebenheit (Kr\"{u}mmungsmass) in einem
gegebenen Punkte und einer gegebenen Fl\"{a}chenrichtung.  Zur
Bestimmung ihrer Massverh\"{a}ltnisse ist es (unter gewissen
Beschr\"{a}nkungen) zul\"{a}ssig und hinreichend, dass die
Kr\"{u}mmungsmass in jedem Punkte in
$\displaystyle n \frac{n - 1}{2}$
Fl\"{a}chenrichtungen beliebig gegeben wird.

\item[\textmd{\S.~3.}]
Geometrische Erl\"{a}uterung.

\item[\textmd{\S.~4.}]
Die ebenen Mannigfaltigkeiten (in denen das Kr\"{u}mmungsmass
allenthalben $= 0$ ist) lassen sich betrachten als einen
besonderen Fall der Mannigfaltigkeiten mit constantem
Kr\"{u}mmungsmass.  Die k\"{o}nnen auch dadurch definirt werden,
dass in ihnen Unabh\"{a}ngigkeit der $n$~fach ausgedehnten
Gr\"{o}ssen vom Ort (Bewegbarkeit derselben ohne Dehnung)
stattfindet.

\item[\textmd{\S.~5.}]
Fl\"{a}chen mit constantem Kr\"{u}mmungsmasse.
\end{description}

\item[\textmd{III.}]
Anwendung auf den Raum.

\begin{description}
\item[\textmd{\S.~1.}]
Systeme von Thatsachen, welche zur Bestimmung der
Massverh\"{a}ltnisse des Raumes, wie die Geometrie sie
voraussetzt, hinreichen.

\item[\textmd{\S.~2.}]
In wie weit ist die G\"{u}ltigkeit dieser empirischen Bestimmungen
wahrscheinlich jenseits der Grenzen der Beobachtung im
Un\-mess\-bar\-grossen?

\item[\textmd{\S.~3.}]
In wie weit im Unendlichkleinen?  Zusammenhang dieser Frage mit
der Naturerkl\"{a}rung.

\end{description}
\end{description}

\end{document}
